\chapter{Meningococcal disease}

\section*{Definition and Overview}
Meningococcal disease is a rare but highly acute and life-threatening bacterial infection caused by the Gram-negative diplococcus \textit{Neisseria meningitidis}. The infection most commonly manifests as meningococcal meningitis (infection of the membranes surrounding the brain and spinal cord) or meningococcal septicaemia (a rapid, invasive bloodstream infection, also known as meningococcemia). Meningococcal disease is characterised by its rapid progression, with the potential to progress from initial non-specific symptoms to septic shock, multi-organ failure, limb gangrene, or death within hours of onset. Prompt clinical recognition and immediate administration of antibiotics are critical to saving lives.

\section*{Epidemiology}
Meningococcal disease occurs globally, with the highest burden of disease found in the sub-Saharan African "meningitis belt." In Australia, the incidence has declined significantly since the introduction of school-aged and infant immunisation programmes, with approximately 100 to 200 cases reported annually. The disease exhibits a seasonal distribution, peaking in late winter and early spring. It predominantly affects children under five years of age (especially infants under twelve months) and adolescents/young adults aged 15 to 24. Of the thirteen identified serogroups, serogroups A, B, C, W, and Y are the primary causes of invasive disease, with serogroup B being the most common cause of cases in Australia.

\section*{Aetiology and Pathophysiology}
The invasive disease arises from the hematogenous spread of the bacteria following nasopharyngeal colonization:
\begin{itemize}
    \item \textbf{Transmission of bacteria:} \textit{Neisseria meningitidis} is transmitted between humans via respiratory droplets or direct contact with oral secretions (e.g., coughing, sneezing, kissing, or sharing drink containers).
    \item \textbf{Nasopharyngeal carriage:} The bacteria colonise the mucosal lining of the nasopharynx; while most carriers remain asymptomatic, the bacteria can occasionally cross the mucosal barrier to enter the bloodstream.
    \item \textbf{Meningococcal septicaemia:} Once in the bloodstream, the bacteria multiply rapidly and shed outer membrane endotoxins, triggering an intense systemic inflammatory response, capillary leakage, vasodilatory shock, and disseminated intravascular coagulation (DIC).
    \item \textbf{Meningitis pathogenesis:} The bacteria cross the blood-brain barrier, multiplying in the subarachnoid space and causing severe purulent inflammation of the meninges, leading to cerebral oedema and elevated intracranial pressure.
\end{itemize}

\section*{Clinical Features}
The clinical presentation of meningococcal disease begins with non-specific, flu-like symptoms before rapidly progressing to characteristic features of meningitis or septicaemia:
\begin{itemize}
    \item \textbf{Early constitutional symptoms:} High fever, severe headache, general malaise, nausea, vomiting, joint pain, and severe muscle pain (especially in the legs).
    \item \textbf{Non-blanching petechial rash:} Small, red or purple spots on the skin that do not fade when pressed with a glass; these can rapidly coalesce into larger purpuric lesions (purpura fulminans).
    \item \textbf{Meningeal irritation signs:} Photophobia (extreme sensitivity to light), neck stiffness, and a positive Kernig's or Brudzinski's sign.
    \item \textbf{Circulatory shock:} Cold hands and feet, rapid breathing (tachypnoea), tachycardia, hypotension, and pale or mottled skin with prolonged capillary refill time.
    \item \textbf{Altered consciousness:} Severe lethargy, confusion, irritability, delirium, or drowsiness that progresses to coma.
    \item \textbf{Paediatric-specific signs:} In infants, typical features include a high-pitched cry, bulging fontanelle, floppy limbs, extreme irritability, and refusal to feed.
\end{itemize}

\section*{Differential Diagnosis}
Because the early features are non-specific, clinicians must differentiate meningococcal disease from other acute febrile illnesses and causes of meningitis:
\begin{itemize}
    \item \textbf{Other bacterial meningitis:} Meningitis caused by \textit{Streptococcus pneumoniae} or \textit{Haemophilus influenzae} type b, which present identically and require microbiological confirmation.
    \item \textbf{Viral meningitis or encephalitis:} Infections caused by enteroviruses or herpes simplex virus, which typically present with milder symptoms and lack the characteristic septic shock or petechial rash.
    \item \textbf{Severe viral infections:} Influenza or severe dengue fever, which can present with high fever, severe myalgia, and rash, but do not cause rapid cardiovascular collapse.
    \item \textbf{Henoch-Sch\"{o}nlein purpura (HSP):} A systemic vasculitis presenting with a purpuric rash (typically on the lower limbs), but without high fever, toxicity, or signs of meningeal irritation.
    \item \textbf{Thrombocytopenic purpura:} Autoimmune or drug-induced thrombocytopenia presenting with petechiae, but lacking the systemic signs of acute bacterial infection.
\end{itemize}

\section*{When to Seek Medical Care}
Meningococcal disease is a medical emergency. Immediate medical attention must be sought at a hospital emergency department if meningococcal disease is suspected.

\textbf{Call 000 (or local emergency services) for:}
\begin{itemize}
    \item \textbf{A non-blanching rash:} Any red or purple spots on the skin that do not fade when pressed with a glass (the "glass test").
    \item \textbf{Signs of meningitis:} Severe headache accompanied by neck stiffness and extreme sensitivity to light.
    \item \textbf{Signs of shock:} Cold hands and feet, rapid breathing, mottled skin, and confusion.
    \item \textbf{An unresponsive infant:} A floppy baby, a high-pitched cry, a bulging soft spot on the head, or extreme listlessness.
\end{itemize}

\section*{Diagnosis and Investigations}
Investigations should be performed urgently, but antibiotic administration must not be delayed for diagnostic procedures:
\begin{itemize}
    \item \textbf{Blood and CSF cultures:} Urgent blood cultures and, if clinically safe, cerebrospinal fluid (CSF) cultures to isolate \textit{Neisseria meningitidis} and perform antibiotic susceptibility testing.
    \item \textbf{Polymerase Chain Reaction (PCR) testing:} Highly sensitive molecular testing of blood or CSF to detect meningococcal DNA, which is particularly useful if antibiotics were administered prior to sampling.
    \item \textbf{Lumbar puncture and CSF analysis:} Examination of CSF showing elevated white blood cells (predominantly neutrophils), low glucose, elevated protein, and Gram-negative diplococci on Gram stain.
    \item \textbf{Full blood count and coagulation profile:} To identify leukocytosis (or leukopenia in severe cases), thrombocytopenia, and prolonged coagulation times (elevated INR and APTT, low fibrinogen) indicating DIC.
\end{itemize}

\section*{Management}
Management requires immediate antibiotic administration, aggressive resuscitation, and critical care support:
\begin{itemize}
    \item \textbf{Empirical antibiotic therapy:} Immediate intravenous (or intramuscular if IV access is delayed) administration of ceftriaxone or benzylpenicillin as soon as the diagnosis is suspected (often prior to hospital transfer).
    \item \textbf{Fluid resuscitation:} Aggressive administration of intravenous crystalloid fluids to treat septic shock and maintain end-organ perfusion.
    \item \textbf{Inotropic and vasoactive support:} Use of vasoactive drugs (such as noradrenaline) in an intensive care setting for patients with fluid-refractory shock.
    \item \textbf{Adjuvant corticosteroids:} Dexamethasone administered just before or with the first dose of antibiotics in cases of suspected meningitis to reduce neurological complications (such as hearing loss).
    \item \textbf{Chemoprophylaxis for close contacts:} Administering oral ciprofloxacin or rifampicin, or intramuscular ceftriaxone, to close household contacts to eradicate nasopharyngeal carriage and prevent secondary cases.
\end{itemize}

\section*{Complications}
The severe inflammatory response and tissue ischaemia associated with the disease can lead to permanent, disabling complications:
\begin{itemize}
    \item \textbf{Neurological sequelae:} Permanent sensorineural hearing loss, intellectual disability, epilepsy, cerebral palsy, or focal motor deficits.
    \item \textbf{Limb amputation:} Extensive tissue necrosis and gangrene of the extremities (purpura fulminans) due to microvascular thrombosis, requiring surgical debridement or amputation.
    \item \textbf{Renal failure:} Acute tubular necrosis and renal failure secondary to prolonged hypotension and microvascular clotting.
    \item \textbf{Adrenal haemorrhage:} Bilateral adrenal gland apoplexy (Waterhouse-Friderichsen syndrome) leading to acute adrenal crisis and profound, refractory shock.
    \item \textbf{Chronic skin scarring:} Extensive skin necrosis requiring skin grafting, resulting in permanent, painful scars.
\end{itemize}

\section*{Prognosis and Outcomes}
The prognosis of meningococcal disease is highly dependent on the speed of diagnosis and initiation of antibiotic therapy. Even with optimal modern intensive care, the overall mortality rate remains at 5\% to 10\%, and can exceed 20\% in patients who present with severe meningococcal septicaemia. Among survivors, approximately 10\% to 20\% will experience permanent, life-altering complications, such as limb amputations, hearing loss, or major neurological deficits. Most patients who recover without complications do so fully over several weeks.

\section*{Prevention and Self-Care}
Vaccination and rapid public health responses are the only effective means of preventing the disease:
\begin{itemize}
    \item \textbf{Meningococcal vaccination:} Adhering to the National Immunisation Program schedule, which includes the Meningococcal ACWY vaccine (for infants and adolescents) and the Meningococcal B vaccine.
    \item \textbf{Adherence to chemoprophylaxis:} Taking prescribed prophylactic antibiotics promptly if identified by public health units as a close contact of a confirmed case.
    \item \textbf{Avoid sharing personal items:} Educating children and adolescents to avoid sharing water bottles, cups, cutlery, or toothbrushes to minimise droplet transmission.
    \item \textbf{Early recognition:} Learning the symptoms of the disease and performing the glass test on any new rash.
\end{itemize}

\section*{Sources and Further Reading}
The following organisations provide detailed, evidence-based resources and information on meningococcal disease:
\begin{itemize}
    \item Australian Government Department of Health and Aged Care. \textit{Meningococcal disease.} https://www.health.gov.au/diseases/meningococcal-disease
    \item Healthdirect Australia. \textit{Meningococcal disease.} https://www.healthdirect.gov.au/meningococcal-disease
    \item NHS. \textit{Meningitis.} https://www.nhs.uk/conditions/meningitis/
    \item Mayo Clinic. \textit{Meningitis - Symptoms and causes.} https://www.mayoclinic.org/diseases-conditions/meningitis/symptoms-causes/syc-20350508
    \item World Health Organization (WHO). \textit{Meningococcal meningitis.} https://www.who.int/news-room/fact-sheets/detail/meningococcal-meningitis
    \item Meningitis Centre Australia. \textit{Signs and Symptoms.} https://meningitis.org.au/signs-and-symptoms/
\end{itemize}
